\documentclass[11pt]{article}

\usepackage[margin=1in]{geometry}
\usepackage{microtype}
\usepackage{amsmath}
\usepackage{graphicx}
\usepackage{booktabs}
\usepackage[hidelinks]{hyperref}

\graphicspath{{figures/}}

\newcommand{\natex}{\texttt{natex}}

\title{A Sharp Bend, No Verdict: Difference-in-Kinks Tests of US Business AI
Adoption at DeepSeek-R1}
\author{Hauke Hillebrandt\thanks{University College London. Email:
\texttt{ucjthhi@ucl.ac.uk}. This paper and the underlying \natex{} software were
prepared with substantial assistance from Anthropic's Claude models; the author
reviewed the analyses and text and is responsible for all remaining errors.}}
\date{July 2026 \\ \small analysis pass of record: 2026-07, \natex{} v0.2.0, seed 0}

\begin{document}

\maketitle

\begin{abstract}
\noindent US Census Bureau Business Trends and Outlook Survey (BTOS) data show a
sharp, well-localized acceleration in business AI use among AI-exposed sectors
(information; finance and insurance; professional/scientific/technical services)
beginning with the first survey wave after the release of DeepSeek-R1 on 20
January 2025. A sector-by-wave difference-in-differences scan localizes the
break at exactly that wave (scan $p=0.010$--$0.030$); the exposed-minus-unexposed
gap slope more than triples, from $+4.2$ to $+13.9$ pp/yr;
and a group difference-in-kinks (DiK) estimate at the release date is nominally
significant ($+7.3$ pp/yr, cluster-robust se $2.8$). Attribution to DeepSeek-R1,
however, fails the falsification battery: placebo cutoffs at non-event 2024
dates produce DiKs as large as
or larger than the R1 estimate ($+11.4$ pp/yr at $t=2024.30$; $+4.0$ at
$t=2024.50$; both significant); only 2 of 6 specification variants at R1 are
significant; equally large local kinks appear in non-exposed sectors; the
predicted \emph{deceleration} at OpenAI's o1 release is insignificant at the
primary bandwidth and sign-flips at the wider one; and the cutoff week is
confounded by the AI diffusion export rule (13 January), the Stargate
announcement (21 January), and a BTOS sample-year rollover in mid-December
2024. The January-2025 acceleration is real but descriptive: no causal
attribution to DeepSeek-R1 survives triage.
\end{abstract}

\section{Introduction}\label{sec:intro}

On 20 January 2025 the Chinese lab DeepSeek released R1, an openly licensed
reasoning model trained at a reported cost far below frontier-lab budgets
\cite{guo2025}. The release was treated by markets as news about the cost curve
of AI: on 27 January 2025 Nvidia alone lost roughly \$590 billion of market
capitalization, the largest one-day loss on record, and an event-study
literature on the ``DeepSeek shock'' emerged almost immediately.\footnote{See,
e.g., ``Nvidia drops \$600B off its market cap amid the rise of DeepSeek,''
TechCrunch, 27 January 2025. We do not rely on this financial-markets
literature for identification; it motivates the candidate event date only.} A natural real-economy question
follows: did cheap, open reasoning models bend the \emph{adoption} curve --- the
share of US businesses actually using AI in production?

Measurement of that adoption curve is recent. The Census Bureau's Business
Trends and Outlook Survey (BTOS) has asked a probability sample of US firms
since September 2023 whether they used AI in producing goods or services in the
last two weeks; \cite{bonney2024} introduce these data and document a rise in
the national use rate from 3.7\% to 5.4\% over the first six months,
concentrated in the Information sector and larger firms. Complementary
measurement includes the 2018 Annual Business Survey module of
\cite{mcelheran2024} (pre-LLM adoption below 6\% of firms, over 18\%
employment-weighted) and the multi-country executive survey of
\cite{yotzov2026} (high headline adoption, small realized effects). This
literature is descriptive by design; formal tests of whether adoption
\emph{trends} broke at specific AI events are, to our knowledge, absent.

A trend break for one group relative to another at a known date is the estimand
of the difference-in-kinks (DiK) design formalized by \cite{bockerman2025},
building on the regression kink design of \cite{card2015}. Kink estimates in
time-series settings are known to produce spuriously significant slope changes,
which is why placebo distributions at non-event dates are the appropriate
yardstick \cite{ganong2018}. This paper applies the \natex{} toolkit's DiK and
scan-based DiD estimators \cite{natex,herlands2018} to the BTOS sector panel
with DeepSeek-R1 as the candidate cutoff, o1 (12 September 2024) as a secondary
contrast, and a mandatory falsification battery, which converts an apparently
significant event-study result into an explicit non-attribution.

\section{Data}\label{sec:data}

The primary source is the Census Bureau's BTOS AI core workbook (DRB approval
CBDRB-FY25-0425): sector-level biweekly estimates and standard errors for
Question 7, ``In the last two weeks, did this business use Artificial
Intelligence (AI) in producing goods or services?'', answer ``Yes,'' in
percentage points. The panel covers 20 sectors (19 two-digit NAICS sectors plus
an unclassified residual) over 54 biweekly waves, 202319--202520 (collection
periods 2023-09-11 through 2025-10-05); 199 of 1{,}080 sector-waves are
disclosure-suppressed, leaving 881 usable cells. The panel predates the
2025-12-03 questionnaire rewording entirely. The running variable $t$ is the
fractional year of each wave's collection-period midpoint. The exposed group is
\{51 information, 52 finance and insurance, 54 professional/scientific/technical
services\}; the kink-leg comparison group is \{11, 21, 23, 72\}. A placebo
outcome (the archived remote-work Question 6 ``Yes'' share) comes from the BTOS
archive workbook (CBDRB-FY23-0478/FY24-0474/FY25-0117), waves 202417--202514.

Coverage is the binding data limitation: sectors 11, 21, 22, and 55 are almost
entirely suppressed (3, 3, 4, and 8 usable waves of 54), so the ``unexposed''
comparison group is effectively sectors 23 (construction) and 72
(accommodation and food services).

\section{Design and Methods}\label{sec:methods}

\paragraph{Group difference-in-kinks.} Following \cite{bockerman2025}, the DiK
estimand is the change in the slope of the outcome in calendar time at a known
cutoff for the exposed group \emph{minus} the same slope change for the
comparison group, estimated by local linear regression on each side of the
cutoff within each group (\natex{} v0.2.0 \cite{natex}). The primary
specification uses cutoff $t_0=2025.055$ (DeepSeek-R1; first post wave 202503,
collected 2025-01-27 to 02-09), bandwidth 0.5 years ($\pm 13$ waves),
triangular kernel, a one-wave donut, and cluster-robust (CR1) standard errors
clustered by sector with 7 clusters and critical values from $t(G-1)$. With 7
clusters and only 2--4 per DiK cell, CR1 inference is fragile even with the
small-sample correction \cite{cameron2015}; we therefore lean on design-based
falsification rather than on the standard errors.

\paragraph{Falsification battery.} (i) a specification grid over bandwidth
$\in \{0.25, 0.5, 1.0\}$ years, kernel $\in$ \{triangular, uniform\}, and donut
$\in \{0, 1\ \text{wave}\}$; (ii) placebo cutoffs at four non-event dates
(2024.30, 2024.50, 2024.85, 2025.30), in the spirit of the permutation logic of
\cite{ganong2018}; (iii) a placebo-in-space battery of per-sector local kinks at
the R1 date; (iv) a placebo \emph{outcome} (remote work) at the same cutoff; and
(v) a secondary contrast at o1 (2024-09-12, $t=2024.6967$): if R1 mattered over
and above the reasoning-model wave o1 began, the exposed-minus-unexposed gap
should \emph{decelerate} there.

\paragraph{Sector-by-wave DiD.} The \natex{} survey pipeline's sector-by-wave
DiD scan (SuDDDS) searches over candidate break dates $t_0$ and window widths
$W$, maximizing a likelihood-ratio scan statistic for a declared-treatment DiD,
calibrated by permutation under an AR(1)-by-unit null \cite{herlands2018}. We
also fit a manual two-way fixed-effects (TWFE) level DiD with wave and sector
fixed effects, CR1 by sector, and a placebo-in-space randomization check.

\paragraph{Honest-inference notes, stated as run.} The pipeline's refusals are
reported, not patched over: the SuDDDS anticipation test refused (Holm
$p=[\mathrm{NaN},\mathrm{NaN},\mathrm{NaN}]$); the dd/synthetic-control/GESS
effect estimators refused with $\tau=\mathrm{NaN}$ (0 usable placebo units on
this unbalanced panel, $\geq 5$ required), hence the manual TWFE estimate
below; and the automated survey's own family verdicts were
discarded as role-assignment artifacts (its ``credible'' RDD, scan $p=0.010$,
$\tau_{\text{2sls}}=0.749$, se $0.343$, is the mechanical rediscovery of
candidate treatment $=$ post with outcome $=t$; its kink family, Holm
$p=0.807$, tested outcome $=$ sector --- vacuous; its DiD family, $p=0.70$,
used a time-invariant treated indicator rather than the declared design;
iv/sc/bunching returned needs-input; dee was null with 1 usable experiment).

\section{Results}\label{sec:results}

\begin{figure}[t]
\centering
\includegraphics[width=\textwidth]{fig1}
\caption{(a) BTOS AI-use share by exposure group across all 54 waves, with the
o1 and DeepSeek-R1 release dates marked. (b) Group difference-in-kinks
estimates (primary specification: bandwidth 0.5 yr, triangular kernel, one-wave
donut) at the R1 cutoff, the o1 cutoff, and four non-event placebo dates, with
95\% CR1 confidence intervals; filled markers are significant at the 5\% level
against $t(G-1)$ critical values. Two non-event placebos are significant, one
larger than the R1 estimate.}
\label{fig:fig1}
\end{figure}

\paragraph{The bend is real \ldots} The SuDDDS scan localizes a break at
$t_0=2025.089$ --- exactly the first post-R1 wave (202503) --- with window
$W=0.536$ yr: max LLR $=22.65$, scan $p=0.030$ on the 20-sector panel, and LLR
$=32.70$, $p=0.010$ (the $q=99$ permutation floor) on the 7-sector subpanel;
the composition check passes ($p=1.000$). The SE-weighted exposed-minus-unexposed
gap slope over $\pm 13$ waves rises from $+4.197$ (se $2.089$) to $+13.940$
(se $1.649$) pp/yr, $\Delta = +9.74$ (se $2.66$). The primary group DiK at R1
(Table~\ref{tab:dik}, first row) is $\hat\tau = +7.297$ pp/yr (CR1 se $2.797$,
95\% CI $[0.454, 14.140]$, $|t|=2.61$ against $t(6)$ critical value $2.447$,
$n=123$): the exposed group's slope jumps from $3.234$ to $12.726$ pp/yr against
$0.485$ to $2.680$ for the comparison group. The manual TWFE level DiD gives
$\hat\tau = +4.583$ pp (CR1 se $0.526$, $t=8.72$), and placebo-in-space is
clean: 0 of 13 control-sector pseudo-treatments reach the minimum exposed
$|t|$ (controls' max $7.76$ vs.\ exposed $10.19/10.49/13.19$; exact-$p$ floor
$1/14 = 0.071$), with per-sector effects of $+4.26$ (51), $+4.14$ (52), and
$+5.35$ (54) pp.

\paragraph{\ldots\ but attribution to R1 fails.} Five findings block a causal
reading. \emph{First}, the placebo-cutoff battery fails: at the non-event dates
2024.30 and 2024.50 the same DiK specification returns $+11.419$ (se $3.556$,
$|t|=3.21$) and $+4.023$ (se $0.889$, $|t|=4.53$) --- both significant, the
first larger than the R1 estimate itself (Table~\ref{tab:dik}, panel B;
Figure~\ref{fig:fig1}b). \emph{Second}, the R1 estimate is significant in only
2 of 6 specification variants. \emph{Third}, the kink is not
exposure-specific: in per-sector local kinks at R1, exposed sector 52 ranks
1/16 ($+15.41$, $t=6.36$), but non-exposed sectors 71 ($+10.46$, $t=4.24$), 53
($+10.00$, $t=3.82$), and 62 ($+7.94$, $t=6.62$) all exceed exposed sectors 51
($+6.56$, $t=1.03$) and 54 ($+6.50$, $t=3.13$) --- the January-2025 kink is
broad-based. \emph{Fourth}, the o1 deceleration contrast is bandwidth-aliased:
$-5.884$ (se $2.909$, $|t|=2.02 <$ crit $2.571$, not significant, though the
sign matches the prediction) at bandwidth 0.5, but $+3.818$ (se $1.417$,
$|t|=2.69$, significant, sign \emph{flipped}) at bandwidth 1.0; the SE-weighted
gap-slope change at o1 is $+0.59$ (se $2.63$) --- the expected deceleration is
not reproduced. \emph{Fifth}, the TWFE level DiD, though clean in space,
violates parallel trends: the pre-R1 gap slope is $+4.20$ pp/yr (se $2.09$), so
the level estimate is descriptive. The placebo \emph{outcome} is consistent
with a broad shock rather than an AI-specific one but is underpowered: the
remote-work DiK at R1 is $+8.821$ (se $5.351$, $|t|=1.65$, not significant, on
a truncated pre-window) --- same sign and magnitude as the AI outcome.

\begin{table}[t]
\centering
\small
\caption{Group difference-in-kinks estimates, BTOS AI-use share (pp/yr).}
\label{tab:dik}
\begin{tabular}{lrrrc}
\toprule
Specification & $\hat\tau$ & se & $|t|$ & sig.\ 5\% \\
\midrule
\multicolumn{5}{l}{\emph{Panel A: DeepSeek-R1 cutoff ($t_0=2025.055$)}}\\
triangular, bw 0.5, donut (primary) & $+7.297$ & $2.797$ & $2.61$ & $\ast$ \\
triangular, bw 0.5, no donut        & $+6.742$ & $1.188$ & $5.68$ & $\ast$ \\
uniform, bw 0.5, donut              & $+4.477$ & $3.079$ & $1.45$ &  \\
triangular, bw 1.0, donut           & $+2.350$ & $2.934$ & $0.80$ &  \\
uniform, bw 1.0, donut              & $+3.941$ & $2.478$ & $1.59$ &  \\
triangular, bw 0.25, donut          & $+10.488$ & $8.130$ & $1.29$ &  \\
\midrule
\multicolumn{5}{l}{\emph{Panel B: placebo and secondary cutoffs (triangular, donut)}}\\
$t_0=2024.30$ (no event), bw 0.5    & $+11.419$ & $3.556$ & $3.21$ & $\ast$ \\
$t_0=2024.50$ (no event), bw 0.5    & $+4.023$ & $0.889$ & $4.53$ & $\ast$ \\
$t_0=2024.85$ (no event), bw 0.5    & $-1.385$ & $3.118$ & $0.44$ &  \\
$t_0=2025.30$ (no event), bw 0.5    & $-1.680$ & $2.293$ & $0.73$ &  \\
o1, $t_0=2024.6967$, bw 0.5         & $-5.884$ & $2.909$ & $2.02$ &  \\
o1, $t_0=2024.6967$, bw 1.0         & $+3.818$ & $1.417$ & $2.69$ & $\ast$ \\
\bottomrule
\end{tabular}

\medskip
\parbox{0.92\textwidth}{\footnotesize Notes: CR1 standard errors clustered by
sector; significance against $t(G-1)$ critical values with $G \in \{5,6,7\}$
clusters in the estimation window (2.447--2.776). Bandwidths in years
($0.5\ \text{yr} = \pm 13$ waves); donut $=$ one wave. Outcome: sector-level
share answering Yes to BTOS Question 7, percentage points.}
\end{table}

\section{Caveats and Conclusion}\label{sec:conclusion}

Three caveats are decisive. First, a calendar-time cutoff absorbs every same-week
event: the Framework for Artificial Intelligence Diffusion export rule
(2025-01-13) and the Stargate announcement (2025-01-21) fall in the same window
as R1 (2025-01-20), so a DiK at $t=2025.055$ cannot distinguish among them ---
and the BTOS sample-year rollover (wave 202426 $\to$ 202501, mid-December
2024) changes respondent-panel composition three to five weeks before the
cutoff. Second, the comparison group is effectively two sectors: NAICS 11, 21,
22, and 55 are almost entirely disclosure-suppressed, and CR1 inference rests
on 7 clusters with only 2--4 clusters per DiK cell, where cluster-robust
standard errors are known to over-reject \cite{cameron2015}. Third, and most
important, the placebo-cutoff battery finds slope breaks as large as the R1
estimate at dates where nothing happened --- exactly the failure mode placebo
distributions exist to catch \cite{ganong2018}: on this panel, the DiK test
cannot tell the R1 date apart from ordinary undulation in a steep adoption
curve.

The verdict is descriptive-only. A sharp, well-localized acceleration in
measured AI adoption in AI-exposed sectors did begin with the first BTOS wave
after 20 January 2025 --- the localization is data-driven and survives
composition and placebo-in-space checks --- but every test that would pin the
bend to DeepSeek-R1 specifically fails: placebo dates match it, most
specifications miss it, unexposed sectors show it too, and the o1 contrast
sign-flips. Something bent in January 2025; these data cannot say it was R1.

\paragraph{Reproducibility.} All estimates were produced with \natex{} v0.2.0
\cite{natex} at seed 0 from the frozen BTOS workbooks (published at
\url{https://www.census.gov/hfp/btos}; not committed to the repository).
Figure~\ref{fig:fig1} regenerates deterministically from the committed
\texttt{figures/make\_fig.py}, which asserts the primary estimate against the
numbers of record before drawing.

\footnotesize
\begin{thebibliography}{9}
\setlength{\itemsep}{2pt}\setlength{\parskip}{0pt}

\bibitem{bockerman2025}
B\"ockerman, P., Jysm\"a, S., and Kanninen, O. (2025).
\emph{Difference-in-Kinks Design}.
IZA Discussion Paper No.~18313.
\url{https://docs.iza.org/dp18313.pdf}

\bibitem{bonney2024}
Bonney, K., Breaux, C., Buffington, C., Dinlersoz, E., Foster, L., Goldschlag,
N., Haltiwanger, J., Kroff, Z., and Savage, K. (2024).
Tracking firm use of AI in real time: A snapshot from the Business Trends and
Outlook Survey.
NBER Working Paper No.~32319.

\bibitem{cameron2015}
Cameron, A.~C., and Miller, D.~L. (2015).
A practitioner's guide to cluster-robust inference.
\emph{Journal of Human Resources}, 50(2), 317--372.

\bibitem{card2015}
Card, D., Lee, D.~S., Pei, Z., and Weber, A. (2015).
Inference on causal effects in a generalized regression kink design.
\emph{Econometrica}, 83(6), 2453--2483.

\bibitem{ganong2018}
Ganong, P., and J\"ager, S. (2018).
A permutation test for the regression kink design.
\emph{Journal of the American Statistical Association}, 113(522), 494--504.

\bibitem{guo2025}
Guo, D., Yang, D., Zhang, H., et al. (2025).
DeepSeek-R1 incentivizes reasoning in LLMs through reinforcement learning.
\emph{Nature}, 645(8081), 633--638.

\bibitem{herlands2018}
Herlands, W., McFowland~III, E., Wilson, A.~G., and Neill, D.~B. (2018).
Automated local regression discontinuity design discovery.
In \emph{Proceedings of the 24th ACM SIGKDD International Conference on
Knowledge Discovery and Data Mining (KDD~'18)}.

\bibitem{mcelheran2024}
McElheran, K., Li, J.~F., Brynjolfsson, E., Kroff, Z., Dinlersoz, E., Foster,
L., and Zolas, N. (2024).
AI adoption in America: Who, what, and where.
\emph{Journal of Economics \& Management Strategy}, 33(2), 375--415.

\bibitem{natex}
Hillebrandt, H. (2026).
\emph{natex: automated natural-experiment discovery and estimation} (version
0.2.0). Software.
\url{https://github.com/HaukeHillebrandt/natex}

\bibitem{yotzov2026}
Yotzov, I., Barrero, J.~M., Bloom, N., Bunn, P., Davis, S.~J., Foster, K.~M.,
Jalca, A., Meyer, B.~H., Mizen, P., Navarrete, M.~A., Smietanka, P., Thwaites,
G., and Wang, B.~Z. (2026).
Firm data on AI.
NBER Working Paper No.~34836.

\end{thebibliography}

\end{document}
